\capitulo{1}{Introducción}

En los últimos años, las tecnologías cloud se han convertido en una pieza fundamental dentro de las infraestructuras empresariales modernas. Las organizaciones necesitan aplicaciones seguras, escalables y accesibles desde cualquier ubicación, permitiendo gestionar información crítica de forma eficiente y reduciendo la complejidad de mantenimiento de infraestructuras físicas.

Uno de los problemas más frecuentes en entornos empresariales es la gestión de incidencias técnicas y operativas. Muchas organizaciones continúan utilizando herramientas poco integradas como hojas de cálculo, cadenas de correos electrónicos o sistemas internos difíciles de mantener. Esto provoca problemas de trazabilidad, retrasos en la resolución de incidencias y dificultades para monitorizar el estado global del sistema.

El presente Trabajo Fin de Grado propone el análisis, diseño e implementación de un sistema empresarial seguro basado en tecnologías cloud utilizando Microsoft Azure. La solución desarrollada permite registrar incidencias empresariales, clasificarlas automáticamente según su prioridad y gestionar la información de forma segura utilizando servicios cloud gestionados.

La arquitectura propuesta utiliza una API desarrollada en Python y Flask desplegada sobre Azure App Service, integrando además servicios como Azure Storage, Azure Key Vault y herramientas de monitorización y automatización como GitHub Actions y SonarQube.

Durante el desarrollo del proyecto se han aplicado metodologías ágiles basadas en sprints, utilizando herramientas profesionales de gestión de tareas y control de versiones para garantizar la trazabilidad del trabajo realizado.

La memoria se estructura en los siguientes capítulos:

\begin{itemize}
    \item Capítulo 1: Introducción y contexto del proyecto.
    \item Capítulo 2: Objetivos generales y específicos.
    \item Capítulo 3: Conceptos teóricos relacionados con cloud computing y arquitecturas empresariales.
    \item Capítulo 4: Técnicas y herramientas utilizadas.
    \item Capítulo 5: Aspectos relevantes del desarrollo.
    \item Capítulo 6: Trabajos relacionados.
    \item Capítulo 7: Conclusiones y líneas futuras.
\end{itemize}

Finalmente, se incluyen diferentes anexos técnicos relacionados con la planificación, diseño, requisitos, documentación técnica y manuales de usuario.